% =====================================================================
\section{Kết luận}
\label{sec:conclusion}

\subsection{Trả lời ba câu hỏi của đề}
\begin{enumerate}
\item \textbf{Kết quả nào phù hợp nhất?} Trên VG $\cap$ COCO, xếp lại top-50 của CLIP bằng scene graph nâng R@1 test từ
39,32 lên 47,95. Trên MKG-W, kênh đếm cạnh tri thức nâng R@1 test từ 24,89 lên 29,82. Trên MMKG tiếng Việt, cùng kênh này,
không chọn lại tham số nào, nâng R@1 từ 23,35 lên 33,92.
\item \textbf{Tri thức quan hệ giúp ở đâu?} Trên VG, mức tăng so với CLIP tăng theo số vật thể trong câu; câu chỉ nêu vật thể,
không có quan hệ, không tăng (−0,80 điểm). Trên 15 câu gián tiếp tự viết của MKG-W, R@1 tăng từ 40 lên 80; câu trực tiếp không
đổi. Trên 18 câu gián tiếp tiếng Việt, số câu đúng ở hạng 1 tăng từ 2 lên 14, nhưng câu trực tiếp giảm từ 9 xuống 4. Thay đồ
thị bằng đồ thị của ảnh khác, hoặc xáo cạnh tri thức, đều làm kết quả thấp hơn CLIP.
\item \textbf{Vì sao mô hình trả về kết quả này?} Lời giải thích là các đại lượng của chính điểm đồ thị; với mô hình chính
của VG, điểm cuối tách chính xác thành tổng các phần. Bỏ 3 phần đứng đầu lời giải thích làm ảnh mất hạng 1 ở 58,1\% câu
val, so với 9,4\% khi bỏ ngẫu nhiên và 34,9\% ở mô hình cấp~1 gốc trộn điểm qua z-score. Lời giải thích này khớp Shapley
tốt nhất trong các cách không ước lượng Shapley, và có GEA xấp xỉ Shapley trên region graph của Visual Genome
(Mục~\ref{sec:xai-compare}).
\end{enumerate}
Hai kết quả âm được ghi nhận: ở mô hình cấp~1, cạnh đóng góp khoảng 0,5 điểm R@1; ở MKG-W, NativE nhận trọng số 0 khi kết
hợp với kênh đếm cạnh.

\subsection{Hạn chế}
\begin{itemize}
\item \textbf{Chỉ 3 seed.} Mọi kiểm định giữa các cấu hình chỉ dựa trên 3 seed mỗi bên.
\item \textbf{Val được dùng nhiều lần} để chọn từng bước của mô hình chính VG, nên số val hơi lạc quan (từ val sang test giảm
0,96 điểm R@1). Test vẫn chỉ chấm một lần.
\item \textbf{Chưa tách được đóng góp của từng loại cạnh} trong mô hình chính VG: bước đồ thị liền gộp bốn thay đổi, và chưa có
ablation chỉ bỏ cạnh kNN.
\item \textbf{Đếm số lượng coi mọi người là một loại.} Câu ``a lone woman'' trong ảnh nhiều người bị coi là sai số lượng.
Ở 8,4\% node câu có số lượng, bảng $g$ làm đổi node ảnh được ghép, có khi sang một node khác nhãn.
\item \textbf{Lỗi parser.} Thư viện FactualSceneGraph 0.7.3 cắt token cuối của câu dài nhất trong mỗi batch. Khoảng 1,8\%
caption bị ảnh hưởng; 0,26\% (train), 0,24\% (val) và 0,31\% (test) mất một phần chữ. Lỗi xảy ra như nhau ở cả ba split.
Nhóm không parse lại dữ liệu, giữ nguyên các số đã báo, và đã sửa lỗi này trong demo.
\item \textbf{Thuộc tính và chữ trong ảnh chưa được dùng}, dù nhiều ca thất bại liên quan tới chúng.
\item \textbf{Giới hạn của xếp lại.} 3,3\% câu test VG có ảnh đúng nằm ngoài top-50 của CLIP và không cứu được bằng xếp lại.
\item \textbf{MKG-W.} Phương pháp năm 2024 không đóng góp vào dòng chính ($\beta = 0$). Đặc trưng ảnh của NativE (383 chiều,
do tác giả trích sẵn) khác CLIP ViT-B/32 của nhóm. Cache Wikidata phải lấy lại ngày 17/11: 441 thực thể mất nhãn tiếng Anh,
số truy vấn val khớp được thực thể đổi từ 704 xuống 703. Bộ câu tự viết chỉ có 25 câu và chỉ lưu số tổng theo nhóm.
\item \textbf{Lời giải thích.} Độ lớn của term không ước lượng được điểm mất: bỏ 3 phần đứng đầu chỉ làm mất khoảng một
phần năm số điểm mà lời giải thích dự đoán, và bộ ba đứng đầu ít khi quyết định thứ hạng. Bỏ dần theo thứ tự của lời giải
thích chỉ cho đúng phản thực tế nhỏ nhất ở 59,7\% câu, so với 90,6\% của occlusion. Các cách giải thích GNN dùng để so sánh
được điều chỉnh cho reranker, nên kết quả của chúng phụ thuộc cách điều chỉnh; validity ngữ nghĩa chấm bằng WordNet, nên chỉ
đo độ hợp lý.
\item \textbf{MMKG tiếng Việt.} Cỡ mẫu nhỏ, chưa có khoảng tin cậy bootstrap; câu soạn tay chỉ được chấm ở $\alpha=0$. Dữ liệu
ưu tiên di tích nổi tiếng; tỉnh lẫn trước và sau đợt sáp nhập năm 2025; thời kỳ chỉ suy theo năm. Linker không chuẩn hoá hai
kiểu đặt dấu thanh (``hòa''/``hoà''), dù cả hai kiểu đều có trong nhãn và câu; ảnh hưởng chưa đo.
\end{itemize}

\subsection{Hướng phát triển}
Các hướng sau đi thẳng từ các ca thất bại đã phân tích, chưa được thử: (1) cho linker đi thêm cạnh lớp con (\texttt{P279}) hoặc
đếm cạnh 2-hop, để sửa các câu gián tiếp tiếng Việt sai vì đáp án thuộc lớp con; (2) chọn $\alpha$ theo từng câu ở kênh đồ thị
tri thức, như trọng số $a$ học theo câu ở VG, vì một $\alpha$ chung không tối ưu cho cả câu trực tiếp lẫn gián tiếp; (3) đưa
thuộc tính vào node ảnh; (4) ablation riêng cho từng loại cạnh của mô hình chính VG; (5) huấn luyện các thành phần của scene
ranker (Phụ lục~\ref{app:scene-ranker}) trên mô hình chính để biết chúng có cộng thêm vào nhau không; (6) xếp lời giải thích
theo điểm mất đo được khi bỏ từng phần, như cột ``bỏ phần này'' của demo, vì occlusion trung thực hơn lời giải thích hiện
tại cả ở deletion lẫn phản thực tế.
